\chapter{Первое задание.}
\section{Сопряжённое пространство. Теоремы Хана-Банаха и Рисса-Фреше.}
\begin{task}[9.1]
    Доказать, что:
    \begin{enumerate}
        \item $\ell_p^* \cong \ell_q, \ p \in (1, +\infty), \ \frac{1}{p} + \frac{1}{q} = 1$;
        \item $\ell_1^* \cong \ell_\infty$;
        \item $c_0^* \cong \ell_1$;
        \item $c^* \cong \ell_1$.
    \end{enumerate}
    Верно ли, что $\ell_\infty^* \cong \ell_1$?
\end{task}
\begin{solution}
    Пусть $p \in (1, \infty)$ и $q$ -- сопряжённый к нему показатель.
    Положим для всякого $g \in \ell^q$ функционал $l_g\colon \ell^p \to \R$ следующим образом:
    \[
        l_g\colon f \mapsto \sum_{n = 1}^{\infty} f(n)g(n).
    \]
    Заметим, что в силу неравенства Гёльдера -- это корректно определённый непрерывный функционал и его норма не больше $\|g\|_q$.

    В силу условий достижения равенства в неравенстве Гёльдера на $f_g \in \ell^p$, определённом как
    \[
        f_g(n) = |g(n)|^{\frac{1}{p - 1}} \cdot \sgn g(n).
    \]
    Тогда
    \[
        \|f_g\|_p^p = \sum_{n = 1}^{\infty} |g(n)|^{\frac{p}{p - 1}} = \sum_{n = 1}^{\infty} |g(n)|^{q} < +\infty.
    \]
    И, очевидно, что
    \[
        |l_g f_g| = \sum_{n = 1}^{\infty} |g(n)|^{\frac{1}{p - 1} + 1} = \sum_{n = 1}^{\infty} |g(n)|^{q} = \|f_g\|_p \|g\|_q.
    \]
    А значит $\ell^q$ изометрически вкладывается подобным образом в $(\ell^p)^*$.

    Пусть теперь $l \in (\ell^p)^*$.
    Пусть $e_n$ -- стандартный базис.
    Определим $g$ как
    \[
        g\colon n \mapsto l(e_n).
    \]
    Заметим, что на $c_{00}$ определённый таким образом $g$ порождает $l_g$, совпадающий с $l\colon$
    \[
        l(f) = \sum_{n = 1}^{N} l(f(n)e_n) = \sum_{n = 1}^{N} f(n)l(e_n) = \sum_{n = 1}^{N} f(n)g(n) = l_g(f).
    \]
    Определим $f^{(m)}$ следующим образом:
    \[
        f^{(m)} = \sgn g \cdot |g|^{q - 1} \cdot \Ind_{n \leq m}.
    \]
    $p$-норма $f^{(m)}$ выражается как
    \[
        \|f^{(m)}\|_p^p = \sum_{n = 1}^{m} |g(n)|^{p(q - 1)} = \sum_{n = 1}^{m} |g(n)|^{q}.
    \]
    А действие $l$ на $f^{(m)}\colon$
    \[
        l f^{(m)} = \sum_{n = 1}^{m} |g(n)|^{q - 1 + 1} = \sum_{n = 1}^{m} |g(n)|^{q}.
    \]
    Тогда, в силу непрерывности функционала:
    \[
        |l f^{(m)}| \leq \|f^{(m)}\|_p \|l\| \leq \|f\|_p \|l\|.
    \]
    Что в точности даёт оценку на $q$-норму $m$-срезки $g$.
    Переход к супремуму по $m$ даёт неравенство:
    \[
        \|g\|_q \leq \|l\| < \infty.
    \]
    Заметим, что $c_{00}$ -- плотно в $\ell^p$, $l|_{c_{00}} = l_g|_{c_{00}}$ и что $l, l_g$ -- непрерывны на $\ell^p$.
    Тогда они совпадают и на замыкании $c_{00}$, то есть на всём пространстве $\ell^p$, а неравенство Гёльдера нам даёт вторую оценку на нормы, которая показывает и изометричность такого вложения.

    Отображения $g \mapsto l_g$ и $l \mapsto g$ являются взаимно обратными и сохраняют нормы, а значит $(\ell^p)^* \cong \ell^q$ (в смысле изометрического изоморфизма).

    Исследуем теперь сопряжённое к $\ell^{1}$.
    Заметим, что для всякого $g \in \ell^\infty$ порождение функционала $l_g\colon \ell^1 \to \R$ аналогическим образом
    \[
        l_{g} \colon f \mapsto \sum_{n = 1}^{\infty} f(n)g(n),
    \]
    даёт нам в силу неравенства Гёльдера непрерывный функционал с нормой не более чем $\|g\|_\infty$.

    С другой стороны, для всякого $e_n \in \ell^1$ -- вектора стандартного базиса -- функционал действует как
    \[
        |l_g(e_n)| = |g(n)| \leq \|l_g\|.
    \]
    Переход к $\sup$ по $n$ даёт вторую оценку на норму функционала:
    \[
        \|g\|_{\infty} = \sup_{n \in \N} |g(n)| \leq \|l_g\|.
    \]
    Таким образом $\|g\|_\infty = \|l_g\|$ и $\ell^\infty$ изометрически вкладывается в $(\ell^1)^*$.

    Пусть теперь $l \in (\ell^1)^*$.
    Положим $g(n) = l(e_n)$.
    Определённая таким образом функция лежит в $\ell^\infty\colon$
    \[
        \sup_{n \in \N} |g(n)| = \sup_{n \in \N} |l(e_n)| \leq \sup_{n \in \N} \|l\|\|e_n\| = \|l\| < +\infty.
    \]
    Очевидно, что для всякой $f \in c_{00}$ действие $l$ и $l_g$ (функционала, порождённого $g$) совпадают.
    В силу плотности $c_{00}$ в $\ell^1$ и непрерывности $l$ и $l_g$ на нём они совпадают и на всём пространстве -- что и даёт обратное изометрическое вложение и показывает $(\ell^1)^* \cong \ell^\infty$.

    Обратное неверно, то есть $(\ell^{\infty})^* \not\cong \ell^1$, так как если сопряжённое -- сепарабельно, то и базовое пространство -- тоже.
    Но $\ell^1$ -- сепарабельно, а $\ell^{\infty}$ -- нет.

    Исследуем теперь сопряжённое к $c_0$.
    Пусть $g \in \ell^1$.
    Определим $l_g\colon c_0 \to \R$ следующим образом:
    \[
        l_g \colon f \mapsto \sum_{n = 1}^{\infty} f(n)g(n).
    \]
    Так как $c_0 \subset \ell^\infty$ то неравенство Гёльдера продолжает спасать жизни и даёт оценку на норму $l_g$ как
    \[
        \|l_g\| \leq \|g\|_1.
    \]

    Пусть теперь $l \in (c_0)^*$.
    Стандартным образом заведём функцию $g$ как
    \[
        g(n) = l(e_n).
    \]
    Покажем, что $g \in \ell^1$.
    Рассмотрим $f^{(m)} \in c_0$ вида
    \[
        f^{(m)}(n) =
        \begin{cases}
            \sgn g(n), & n \leq m, \\
            0, & n > m.
        \end{cases}
    \]
    Рассмотрим действие $l$ на $f^{(m)}\colon$
    \[
        |lf^{(m)}| = \left|\sum_{n = 1}^{m} \sgn g(n)l(e_n) \right| = \sum_{n = 1}^{m} |g(n)|.
    \]
    Так как
    \[
        |l f^{(m)}| \leq \|l\|\|f^{(m)}\|,
    \]
    и $\|f^{(m)}\| = 1$, то
    \[
        \sum_{n = 1}^{m} |g(n)| \leq \|l\|.
    \]
    А значит и
    \[
        \|g\|_1 \leq \|l\|.
    \]
    Очевидно, что на $c_{00}$ функционалы $l$ и $l_g$ совпадают.
    В силу плотности $c_{00}$ в $c_0$ (так как последовательности убывают на бесконечности, то срезками можно приблизить сколь угодно точно) функционалы $l$ и $l_g$ совпадают и на всём $c_0$.
    Из уже показанных оценок видно, что $\|l\| = \|l_g\| = \|g\|$, что и даёт искомый изометрический изоморфизм $(c_0)^* \cong \ell^1$.

    Рассмотрим теперь $c$ и найдём сопряжённое к нему.
    Стандартный базис для $c$ -- это $e_0 = (1, 1, \ldots)$ и $\{e_n\}_{n = 1}^{\infty}$, где $e_n(m) = \delta_{mn}$.
    Для $x \in c$ запишем разложение по базису:
    \[
        x = \lim_{n \rightarrow \infty} x(n) e_0 + \sum_{k = 1}^{\infty}(x(k) - \lim x(k))e_k.
    \]
    Пусть $f \in c^*$.
    Обозначим $x(\infty) = \lim x(k)$.
    Тогда
    \[
        f(x) = x(\infty)f(e_0) + \sum_{k = 1}^{\infty} (x(k) - x(\infty))f(e_k).
    \]
    Заведём вектор $y = (f(e_0), f(e_1), \ldots)$.
    Покажем, что он лежит в $\ell^1$.
    Рассмотрим $x^{m}$, определяемый как
    \[
        x^{m}(n) = \begin{cases}
                       \sgn y_n, & n \leq m, \\
                       0, & n > m.
        \end{cases}
    \]
    На таком векторе
    \[
        f(x^m) = \sum_{k = 1}^{m} |y_k|,
    \]
    а так как $\|x^m\| = 1$ и $f$ -- непрерывен, то $y \in \ell^1$.

    Чуть по-другому перепишем разложение $f(x)$ по базису:
    \[
        f(x) = x(\infty)\left(f(e_0) - \sum_{k = 1}^{\infty} f(e_k)\right) + \sum_{k = 1}^{\infty} x(k)f(e_k).
    \]
    Пусть $\tilde{y} = \left(f(e_0) - \sum_{k = 1}^{\infty} f(e_k), f(e_1), \ldots\right)$.
    Покажем, что $\|f\| = \|\tilde{y}\|$.
    Рассмотрим $\tilde{x}^{m}$, определяемый как
    \[
        \tilde{x}^{m}(n) =
        \begin{cases}
            \sgn \tilde{y}(n), & n \leq m, \\
            \sgn \tilde{y}(0), & n > m.
        \end{cases}
    \]
    Распишем действие $f$ на $\tilde{x}^{m}$:
    \[
        f(\tilde{x}^{m}) = |\tilde{y}(0)| + \sum_{k = 1}^{m} |\tilde{y}(k)| + \sgn \tilde{y}(0) \sum_{k = m + 1}^{\infty} \tilde{y}(k).
    \]
    Тогда
    \[
        f(\tilde{x}^{m}) \rightarrow |\tilde{y}(0)| + \sum_{k = 1}^{\infty} |\tilde{y}(k)|.
    \]
    Так как $\|\tilde{x}^{m}\| \leq 1$, то $\|f\| \geq \|\tilde{y}\|$.
    Обратное неравенство очевидно получается из неравенство Гёльдера.
    Таким образом, показано что для всякого $f \in c^*$ существует $\tilde{y} \in \ell^1$ такой, что $f(x) = \langle \tilde{y}, x \rangle$ и $\|f\| = \|\tilde{y}\|$.
    Показано и обратное: всякая $\tilde{y} \in \ell^1$ порождает непрерывный функционал с нормой, равной нормой $\tilde{y}$.
    А значит $c^* \cong \ell^1$.
\end{solution}
\begin{task}[9.3]
    Пусть $L$ -- замкнутое линейное подпространство нормированного пространства $X$ и $y \in X, \ y \notin L$.
    Доказать, что найдётся функционал $f$ на $X$ такой, что $f|_{L} \equiv 0, \ f(y) = 1, \ \|f\| = \frac{1}{\rho(y, L)}$.
\end{task}
\begin{solution}
    Определим функционал $\tilde{f}\colon L \oplus \Lin\{y\} \to \Cm$ как $\tilde{f}|_L = 0$ и $\tilde{f}(y) = 1$, продолженное по линейности.
    Его норма равна
    \[
        \|\tilde{f}\| = \sup_{x \neq 0} \frac{|\tilde{f}(x)|}{\|x\|} = \sup_{t \neq 0, z \in L} \frac{|\tilde{f}(ty + z)|}{\|z + ty\|} = \frac{|\tilde{f}(y)|}{\inf_{z \in L} \|y - z\|} = \frac{1}{\rho(y, L)}.
    \]
    Рассмотрим мажоранту $p\colon X \to \R$, имеющую вид
    \[
        p(x) = \|\tilde{f}\|\|x\|.
    \]
    Очевидно, что $|\tilde{f}(x)| \leq \|\tilde{f}\|\|x\|$ на $L \, \oplus \, \Lin\{y\}$.
    Тогда, согласно теореме Хана-Банаха, существует $f$ -- продолжение на всё $X$ с сохранением мажоранты.
    В силу сохранения мажоранты $\|f\| \leq \frac{1}{\rho(y, L)}$.
    А поскольку $f$ является продолжением $\tilde{f}$, то норма в точности равна $\frac{1}{\rho(y, L)}$.
\end{solution}
\begin{task}[9.11]
    Найти норму функционала $\phi$ на пространстве $C[a, b]\colon$
    \[
        \phi(f) = \int_a^b f(x)g(x)dx,
    \]
    где $g$ -- фиксированная непрерывная функция.
    Исследовать вопрос о том, когда норма достигается.
\end{task}
\begin{solution}
    Непрерывная функция достигает своего супремума на $[a, b]$ и, следовательно
    \[
        \forall f \in C[a, b]\colon \|f\|_{\infty} \leq 1 \hookrightarrow |\phi(f)| \leq \|f\|_{\infty}\|g\|_{1} \leq \|g\|_{1}.
    \]
    Значит $\|\phi\| \leq \|g\|_{1}$.

    Рассмотрим функции вида
    \[
        f_n = \frac{g}{\sqrt{g^2 + \frac{1}{n^2}}}.
    \]
    Заметим, что $\|f_n\|_\infty \leq 1$.
    Тогда
    \[
        |\phi(f_n)| = \left|\int_a^{b} \frac{g^2(x)}{\sqrt{g^2(x) + \frac{1}{n^2}}}dx\right| = \int_a^b \frac{g^2(x)}{\sqrt{g^2(x) + \frac{1}{n^2}}}dx \rightarrow \|g\|_1,
    \]
    а значит норма $\phi$ в точности равна $\|g\|_1$.

    Пусть $Z(g) = g^{-1}(\{0\})$.
    Покажем, что норма достигается тогда и только тогда, когда $Z(g)$ не содержит изолированных точек изменения знака $\sgn$.

    Очевидно, что если $Z(g)$ не содержит изолированных точек смены знака, то существует непрерывная (кусочно-линейная, например) функция $h$ такая, что
    \[
        h|_{[a, b] \setminus Z(g)} = \sgn g|_{[a, b] \setminus Z(g)}.
    \]
    Но тогда
    \[
        \int_{a}^{b} h(x)g(x)dx = \int_{Z(g)} g(x)dx + \int_{[a, b] \setminus Z(g)} |g(x)|dx = \|g\|_1.
    \]

    Если же норма достигается на некоторой $f\colon \|f\|_\infty \leq 1$, то стоит заметить, что
    \[
        \int_{a}^{b}|g(x)|dx = \left|\int_a^b f(x)g(x)dx\right| \leq \int_{a}^{b} |f(x)||g(x)|dx,
    \]
    а значит
    \[
        \int_{a}^{b} (1 - |f(x)|)|g(x)|dx \leq 0.
    \]
    Но подынтегральная функция неотрицательна, а следовательно она равна $0$ почти всюду.
    В частности, $|f(x)| = 1$ на $[a, b] \setminus Z(g)$ и $f(x) = \sgn g(x)$ на $[a, b] \setminus Z(g)$.
    Но в окрестности изолированной точки изменения знака $\sgn$ не существует непрерывного продолжение $\sgn g(x)$ на $Z(g)$, а следовательно $f$ не может быть непрерывной.
\end{solution}
\begin{task}[9.4]
    Привести пример функционала в пространстве $C[a, b]$, не достигающего своей нормы.
\end{task}
\begin{solution}
    Пример очевидно строится по задаче 9.11.
    Так, нам достаточно взять в качестве ядра $g(x) = \frac{x - a}{b - a} - \frac{1}{2}$ и в таком виде $g(a) = -\frac{1}{2}$, $g(b) = \frac{1}{2}$, имеется единственный $0$, который изолирован.
    Определённый с таким ядром интегральный оператор не будет достигать своей нормы по критерию из 9.11.
\end{solution}
\begin{task}[9.5]
    Доказать, что непрерывный линейный функционал $f$ в нормированном пространстве $X$ достигает своей нормы тогда и только тогда, когда для некоторого (и тогда для любого) элемента $x \in X \setminus \ker f$ существует элемент наилучшего приближения в $\ker f$.
\end{task}
\begin{solution}
    Случай нулевого функционала тривиален.
    Пусть $f \neq 0$.
    Существует $z\neq 0\colon f(z) \neq 0$.
    Так как $X = \ker f \oplus \Lin\{z\}$, то всякий вектор представляется единственным образом как $x = y + tz$.
    Тогда
    \[
        \|f\| = \sup_{y \in \ker f, t > 0} \frac{|f(y + tz)|}{\|y + tz\|} = \sup_{y \in \ker f, t > 0} \frac{t|f(z)|}{\|y + tz\|} = \sup_{y \in \ker f} \frac{|f(z)|}{\|y + z\|} = \frac{|f(z)|}{\rho(z, \ker f)}.
    \]
    Видно, что существование элемента наилучшего приближения $z$ в $\ker f$ эквивалентно существованию вектора, на котором достигается норма функционала -- что и требовалось показать.
\end{solution}
\begin{task}[9.7]
    Пусть $E$ -- банахово, $\{x_n\}_{n = 1}^{\infty} \subset E$ и $\sup_{n} |f(x_n)| < \infty$ при $f \in E^*$.
    Доказать, что $\sup \|x_n\| < \infty$.
\end{task}
\begin{solution}
    Рассмотрим семейство операторов $T_n\colon$
    \[
        T_n \colon f \mapsto f(x_n).
    \]
    Заметим, что $\{T_n\}$ -- поточечно ограничено по условию, то есть
    \[
        \forall f \in E^* \hookrightarrow \sup_{n \in \N} |T_n f| = \sup_{n \in \N} |f(x_n)| < \infty.
    \]
    Следовательно, в силу банаховости $E^*$ и теоремы Банаха-Штейнгауза, семейство $\{T_n\}$ является равномерно ограниченным:
    \[
        \sup_{n \in \N} \|T_n \| < \infty.
    \]
    Но
    \[
        \|T_n\| = \sup_{f \in E^*, \|f\| = 1} \|T_n f \| = \sup_{f \in E^*, \|f\| = 1} |f(x_n)| = \|x_n\|,
    \]
    где последнее равенство берётся из того, что $|f(x_n)| \leq \|f\|\|x_n\| = \|x_n\|$ и если мы рассмотрим линейный функционал $g\colon \Lin\{x_n\} \to \Cm, \  \alpha x_n \mapsto \alpha \|x_n\|$, то его норма на $\Lin\{x_n\}$ единична, и существует продолжение с её сохранением на всё $E$ (а значит $\|T_n\| \geq \|x_n\|$ -- что и показывает искомое равенство).

    Тогда и семейство $\{\|x_n\|\}_{n \in \N}$ является равномерно ограниченным.
\end{solution}
\begin{task}[9.9]
    Пусть $L \subset H$ -- линейное многообразие в гильбертовом пространстве, а $f$ -- непрерывный линейный функционал на $L$.
    Доказать, что существует единственный $\tilde{f} \in H^*\colon \tilde{f}|_{L} = f$ и $\|\tilde{f}\| = \|f\|$.
\end{task}
\begin{solution}
    Существует единственное продолжение $f$ на $\overline{L}$ (по непрерывности).
    Оно сохраняет норму.
    Далее будем считать $L$ -- замкнутым подпространством $H$.
    Замкнутое подпространство $H$ само по себе является гильбертовым пространством и, следовательно, согласно теореме Рисса-Фреше существует единственный $y \in L\colon$
    \[
        \forall x \in L\hookrightarrow \langle f, x \rangle = \langle x, y \rangle,
    \]
    а $\|y\| = \|f\|$.
    Определим продолжение $\tilde{f}$ на всё $H$ как
    \[
        \forall x \in H \hookrightarrow \langle \tilde{f}, x \rangle = \langle x, y \rangle.
    \]
    Его норма, очевидно, совпадает с нормой $f$.

    Пусть существует несколько различных продолжений $f_1$ и $f_2$.
    В силу теоремы Рисса-Фреше им соответствуют векторы $y_1$ и $y_2$.
    Поскольку проекции на $L$ у них должны совпадать (в силу равенства норм), то нормы остатков равны $0$ и, следовательно, проекции на $L^{\perp}$ равны $0$
\end{solution}
\section{Слабая и слабая-* топология.}
\begin{task}[10.3]
    Пусть $f_n(x) = \sin nx$.
    Показать, что $f_n$ в $L^2([-\pi, \pi])$ сходится слабо, но не сильно.
\end{task}
\begin{solution}
    Заметим, что в силу того, что коэффициенты ряда Фурье убывают по модулю на бесконечности, то $\forall f \in L^2([-\pi, \pi])$ верно, что
    \[
        \langle f, \sin nx \rangle \rightarrow 0, n \rightarrow \infty.
    \]
    Тогда $f_n \xrightarrow{w} 0, n \rightarrow \infty$ (согласно критерию слабой сходимости).

    С другой стороны, $\|f_n\| = \sqrt{\pi}$ и, следовательно, сходимости по норме к $0$ нет -- что и требовалось показать.
\end{solution}
\begin{task}[10.4]
    Пусть множество $M \subset L^2([-\pi, \pi])$ состоит из функций вида $f_{m, n} = \sin mx + m \sin nx$.
    Показать, что первое слабое секвенциальное замыкание $M$ не совпадает со вторым.
\end{task}
\begin{solution}
    Очевидно, что $0 \in \left[\left[ M \right]^{w, seq}\right]^{w, seq}$, так как
    \[
        f_{m, n} \xrightarrow{w} \sin mx, n \rightarrow \infty.
    \]
    и
    \[
        \sin mx \xrightarrow{w} 0, m \rightarrow \infty.
    \]
    С другой стороны, предположим что $0 \in \left[ M \right]^{w, seq}$.
    Тогда существует последовательность точек $M$ такая, что
    \[
        f_{m_k, n_k} \xrightarrow{w} 0, k \rightarrow \infty.
    \]
    Заметим, что $\{m_k\}$ не может быть неограничена -- иначе последовательность норм $\{\|f_{m_k, n_k}\|\}$ не является ограниченной.
    Значит $m_k$ имеет стационарную подпоследовательность, стабилизирующуюся на некотором $m \in \N$ и
    \[
        f_{m, n_k} \xrightarrow{w} 0, k \rightarrow \infty.
    \]
    Согласно предположению о слабой сходимости к $0$ для $\sin mx\colon$
    \[
        \langle f_{m, n_k}, \sin mx \rangle \rightarrow 0.
    \]
    Но
    \[
        |\langle f_{m, n_k}, \sin mx \rangle| \geq \pi,
    \]
    и таким образом изначальное предположение неверно и $0 \notin [M]^{w, seq}$.
\end{solution}
\begin{task}[10.6]
    Доказать, что в конечномерных нормированных пространствах слабая сходимость совпадает со сходимостью по норме.
\end{task}
\begin{solution}
    Пусть $X$ -- конечномерное нормированное пространство и $\dim X = n$.
    Заметим, что слабая топология всегда вложена в сильную, поскольку предбазой являются множества вида:
    \[
        \sigma_w = \left\{ f^{-1}(U) \mid U \subset \tau_{\Cm}, \ f \in X^* \right\},
    \]
    а прообраз открытого $U$ под действием непрерывного функционала $f$ открыт в сильной топологии (то есть $\sigma_w \subset \tau_{\|\cdot\|}$, а значит и $\tau_w \subset \tau_{\|\cdot\|}$).

    Пусть $B_\epsilon^{\infty}(x) \subset X$ -- открытый $\epsilon$-шар в $\infty$-норме.
    Рассмотрим координатные функционалы $e_i$.
    В слабой топологии множество $V_\delta = \bigcap_{i = 1}^{n} e_i^{-1}(B_{\epsilon - \|y - x\|_\infty}(y_i))$ открыто.
    Заметим, что $V_{\delta} \subset B_{\epsilon}^{\infty}(x)$ и, следовательно, $B_{\epsilon}^{\infty}(x)$ -- открыто в слабой топологии, а значит $\tau_w = \tau_{\|\cdot\|_{\infty}}$.
    Поскольку все нормы на конечномерных пространствах эквивалентны, то слабая топология совпадает с топологией всякой нормы, а значит и слабая сходимость эквивалентна сходимости по норме.
\end{solution}
\begin{task}[10.7]
    Доказать, что из слабой сходимости последовательности элементов $\ell^1$ следует её сходимость по норме.
\end{task}
\begin{solution}
    Заметим, что достаточно показать, что если $x_n \subset \ell^1$ такова, что $x_n \xrightarrow{w} 0$, то и $\|x_n\| \rightarrow 0$.
    И действительно: если $x_n \xrightarrow{w} x$, то это эквивалентно тому, что $\forall f \in (\ell^1)^* \hookrightarrow f(x_n) \rightarrow f(x) \Longleftrightarrow f(x_n - x) \rightarrow 0 \Longleftrightarrow x_n - x \xrightarrow{w} 0$.

    Известно, что $(\ell^1)^* \cong \ell^\infty$.
    Всякий координатный функционал $f_m(k) = \delta_{mk}$ лежит в $\ell^{\infty}$ и
    \[
        f_m x_n = x_n(m) \rightarrow 0.
    \]
    А значит всякая координата последовательности стремится к $0$ при $n \rightarrow \infty$.

    Также известно, что всякая слабо сходящаяся последовательность ограничена по норме.
    Пусть $M = \sup \|x_n\|_1 < \infty$.

    Предположим, что $\|x_n\|_1 \not\rightarrow 0, n \rightarrow \infty$.
    Тогда $\exists \epsilon > 0$ и подпоследовательность $\{x_{n_k}\}$ такая, что $\|x_{n_k}\|_1 \geq \epsilon$.
    Обозначим последовательность $\{x_{n_k}\}$ за $y_n$.
    Заметим, что для всякого $n \in \N$ существует конечное множество индексов $F_n \subset \N$ такое, что
    \[
        \sum_{k \in F_n} |y_n(k)| \geq \frac{\epsilon}{2}.
    \]
    Построим попарно непересекающиеся конечные множества $E_j$ и подпоследовательность $\{y_{n_j}\}$ так, чтобы
    \[
        \sum_{k \in E_j} |y_{n_j}(k)| \geq \frac{\epsilon}{4}.
    \]
    Это возможно, поскольку для конечного множества $W \subset \N$ в силу стремления координат к 0 мы можем выбрать достаточно большое $n$ так, чтобы $\sum_{k \in W} |y_{n}(k)|$ была сколь угодно мала.
    А значит, строя $E_j$ индуктивно, мы на каждом шаге выбирать $W = \bigcup_{i < j} E_i$ и $n_j$ такое, чтобы
    \[
        \sum_{k \in W} |y_{n_j}(k)| \leq \frac{\epsilon}{8}.
    \]
    Тогда
    \[
        \sum_{k \in F_{n_j} \setminus W} |y_{n_j}(k)| \geq \frac{3\epsilon}{8}.
    \]
    Чего мы и хотели.
    Определим теперь функционал $g \in \ell^\infty$ как
    \[
        g(k) = \begin{cases}
                   e^{-i\arg y_{n_j}(k)}, & \exists j\colon k \in E_j \\
                   0.
        \end{cases}
    \]
    Так как $|g| \leq 1$, то $g$ действительно существенно ограничен и лежит в $\ell^{\infty}$.
    Но для $y_{n_j}\colon$
    \[
        g y_{n_j} = \sum_{k} g(k) y_{n_j}(k) \geq \sum_{k \in E_j} |y_{n_j}(k)| - \sum_{k \in E_i, i \neq j} |y_{n_j}(k)| \geq \frac{3\epsilon}{8} - \frac{\epsilon}{4} = \frac{\epsilon}{8}.
    \]
    А значит слабой сходимости к $0$ нет -- противоречие.
\end{solution}
\begin{task}[10.14]
    Пусть $U = \{f \in L^1([0, 1]) \mid |f(t)| \leq 1 \ \text{п.в. } t \in [0, 1] \}$.
    Показать, что $U$ слабо секвенциально компактно в $L^1([0, 1])$.
\end{task}
\begin{solution}
    Заметим, что множество $U \subset L^2([0, 1])$ (в силу неравенства Гёльдера).
    Оно, очевидно, является выпуклым, а значит слабая (и сильная) замкнутость для него эквивалентны (по теореме Мазура).
    Предположим, что $\{x_n\} \subset U$ и $x_n \xrightarrow{\|\cdot\|_2} x_0 \in L^2([0, 1])$.
    Тогда $x_n \xrightarrow{\LL^1} x_0, n \rightarrow \infty$ и по теореме Рисса существует $x_{n_k} \xrightarrow{a.s.} x_0, k \rightarrow \infty$.
    Следовательно $a.s.-\forall t \in [0, 1] \hookrightarrow x_0(t) \leq 1$ и $x_0 \in U$.

    Поскольку $L^2([0, 1])$ -- рефлексивно и $U$ ограничено, то для всякой последовательности $\{x_n\} \subset U$ существует (по теореме Банаха-Тихонова) подпоследовательность $\{x_{n_k}\}$, слабо сходящяся к $x_0$.
    В силу замкнутости $U$ верно, что $x_0 \in U$.
    Заметим, что по критерию слабой сходимости:
    \[
        x_{n_k} \xrightarrow{w} x_0 \Longleftrightarrow \forall y \in L^2([0, 1]) \hookrightarrow \langle x_{n_k}, y \rangle \rightarrow \langle x_0, y \rangle.
    \]
    Но всякий $y \in L^{\infty}([0, 1]) \cong L^1([0, 1])^*$ лежит и в $L^2([0, 1])$, а значит $x_{n_k} \xrightarrow{w} x_0$ и в $L^1([0, 1])$.
    Таким образом $U$ -- слабо секвенциально компактно в $L^1([0, 1])$.
\end{solution}
\begin{task}[T1]
    Доказать, что конечномерное нормированное пространство рефлексивно.
\end{task}
\begin{solution}
    Для доказательства рефлексивности конечномерного нормированного пространства $X$ достаточно показать сюръективность изометрии Банаха $\Phi\colon X \to X^{**}$.
    Так как она инъективна, то
    \[
        \dim \im \Phi = \dim X - \dim \ker \Phi = \dim X.
    \]
    С другой стороны, из линейной алгебры известно, что в конечномерном случае размерность сопряжённого пространства совпадает с размерностью исходного.
    Но, тогда, и для второго сопряжённого верно
    \[
        \dim X^{**} = \dim X^* = \dim X.
    \]
    А значит
    \[
        \dim \im \Phi = \dim X^{**},
    \]
    и изометрия Банаха сюръективна.
\end{solution}
\begin{task}[T2]
    Доказать, что в бесконечномерном нормированном пространстве слабое замыкание единичной сферы -- единичный замкнутый шар.
\end{task}
\begin{solution}
    Заметим, что всякая слабая окрестность содержит бесконечномерное аффинное подпространство.
    Следовательно, никакое слабо открытое множество не может быть ограниченным.
    Таким образом, если $x \in B$, то всякая его слабая окрестность пересекается с единичной сферой $S$ (можно, например, рассмотреть отрезок от $x$ до точки с достаточно большой нормой и использовать теорему о промежуточном значении для нормы).
    А значит слабое замыкание $S$ содержит $B$.
    В то же время, по теореме Мазура, единичный замкнутый шар является слабо замкнутым -- а значит слабое замыкание единичной сферы лежит в нём.
    Что и означает совпадение слабого замыкание сферы с замкнутым единичным шаром.
\end{solution}
\begin{task}[T3]
    Доказать, что если в нормированном пространстве существует последовательность точек единичной сферы $S$, слабо сходящаяся к нулю, то слабое секвенциальное замыкание $S$ совпадает с единичным замкнутым шаром $B$.
\end{task}
\begin{solution}
    Пусть $x_n \xrightarrow{w} 0$, где $\{x_n\} \subset S$.
    Так как слабое секвенциальное замыкание вложено в слабое замыкание, то $[S]^{w, seq} \subset [S]^{w} = B$.

    Покажем обратное включение.
    Пусть $x \in B \setminus S$.
    Для всякого $n \in \N$ рассмотрим функцию
    \[
        g\colon t \mapsto \|x + t x_n\|.
    \]
    Она непрерывна по норме, и $g(t) \rightarrow \infty, t \rightarrow \infty$, а $g(0) = \|x\| < 1$.
    Тогда по теореме о промежуточном значении существует $t_n$ такой, что $g(t_n) = 1$.
    Тогда последовательность $\{x + x_n t_n\}$ лежит на единичной сфере.
    Заметим, что последовательность $\{t_n\}$ -- ограничена (в силу неравенства треугольника $t_n \in [-1 - \|x\|, \|x\| + 1]$).
    Покажем, что $x + x_n t_n \xrightarrow{w} x$.
    Если $f \in X^*$, то
    \[
        f(x + x_n t_n) = f(x) + t_n f(x_n).
    \]
    Так как $f(x_n) \rightarrow f(0) = 0$, а $t_n$ -- ограничена, то $t_n f(x_n) \rightarrow 0$ и
    \[
        f(x + x_n t_n) \rightarrow f(x),
    \]
    и чего в силу произвольного выбора $f$ и следует слабая сходимость к $x$.
    Точка $x \in B \setminus S$ выбрана произвольно, а значит $[S]^{w, seq} \supset B$ и таким образом утверждение задачи показано.
\end{solution}
\begin{task}[T4]
    Пусть $X$ -- бесконечномерное нормированное пространство.
    Показать, что слабая топология в $X$ неметризуема.
\end{task}
\begin{solution}
    Докажем более общее утверждение.
    Пусть $X$ -- бесконечномерное нормированное пространство, $\FF \subset X^*$ -- семейство функционалов, разделяющее точки на $X$ и являющееся бесконечномерным замкнутым подпространством $X^*$.
    Тогда $\FF$-слабая топология $\tau_{\FF}$ на $X$ неметризуема.

    Предположим противное.
    Тогда существует счётная локальная база в $0$ -- некоторое семейство $\{U_n\}$.
    Без ограничения общности мы можем считать, что она состоит из множеств из базы $\FF$-слабой топологии, то есть
    \[
        U_n = \bigcap_{k = 1}^{N_n} V(0, f_k^n, \epsilon_k^n).
    \]
    Любую локальную базу можно сделать убывающей, рассмотрев множества вида $W_N = \bigcap_{n = 1}^{N} U_n$.
    Если теперь её перенумеровать и <<распрямить>>, то мы можем считать что локальная база $\{W_N\}$ имеет вид
    \[
        W_N = \bigcap_{k = 1}^{N} V(0, f_k, \epsilon_k).
    \]
    Рассмотрим множества $E_N \subset \FF$, определяемые как
    \[
        E_N= \Lin\{f_1, \ldots, f_N\},
    \]
    Всякое $E_N$ -- конечномерно, а, следовательно, замкнуто и нигде не плотно в $\FF$.
    Покажем, что тогда
    \[
        \FF = \bigcup_{N = 1}^{\infty} E_N.
    \]
    Пусть $f \in \FF$.
    Множество $V = V(0, f, 1)$ является $\FF$-слабой окрестностью нуля, а значит существует $n \in \N$ такое, что $W_n \subset V$.
    Пусть $f \notin E_n$.
    Тогда существует такой $x$, что $x \in \ker f_j$ при $j = \overline{1, N}$, но $x \notin \ker f$.
    Но, тогда, для всякого $t > 0$ вектор $tx \in \ker f_j$ (и $tx \notin \ker f$).
    Тогда $\Lin\{x\} \subset W_n$.
    С другой стороны, при достаточно большом $t$ число $|f(tx)|$ станет больше единицы и, следовательно, $tx$ не будет лежать в $V$ -- противоречие.
    Значит $f \in E_n$ и $\FF = \bigcup_{n = 1}^{\infty} E_n$.
    Но $\FF$ -- полно (как замкнутое подпространство банахова пространства $X^*$) и представление в виде объединения нигде не плотных множеств противоречит теореме Бэра.
    Таким образом изначальный посыл был неверен и $\FF$-слабая топология не метризуема.

    Слабая топология на $X$ подходит под условия вышесформулированного утверждения и, следовательно, неметризуема.
\end{solution}
\begin{task}[T5]
    Пусть $X$ -- бесконечномерное банахово пространство.
    Показать, что слабая-* топология в $X^*$ неметризуема.
\end{task}
\begin{solution}
    *-слабая топология порождается $X$, вложенными в $X^{**}$.
    Так как $X$ -- бесконечномерное банахово пространство, то она подпадает под условия применения утверждения, сформулированного в предыдущей задаче и, следовательно, не является метризуемой.
\end{solution}
\begin{task}[T6]
    Исследовать последовательность $\{x_n\} \subset l_\infty$ вида
    \[
        x_n(k) = \begin{cases}
                     \cos k, & k \leq n, \\
                     \sin k, & k > n.
        \end{cases}
    \]
    на слабую сходимость в $l_\infty$.
    Рассматривая $l_\infty$ как сопряжённое к $l_1$ исследовать в нём слабую-* сходимость последовательности $x_n$.
\end{task}
\begin{solution}
    Исследуем *-слабую сходимость последовательности.
    Рассмотрим $x\colon k \mapsto \cos k$.
    Пусть $y \in \ell^1$.
    \[
        |y(x_n - x)| = \left|\sum_{k = 1}^{\infty} (x_n(k) - x)y(k)\right| = \left|\sum_{k = n + 1}^{\infty} (\sin k - \cos k) y(k) \right| \leq 2 \sum_{k = n + 1}^{\infty} |y(k)| \rightarrow 0, n \rightarrow \infty.
    \]
    Что и доказывает *-слабую сходимость.

    Рассмотрим координатный функционал $\delta_k\colon \ell^{\infty} \to \R$, который действует как
    \[
        \delta_k \colon x \mapsto x(k).
    \]
    Если $x_n$ слабо сходится, то
    \[
        \delta_k(x_n) \rightarrow \delta_k(x).
    \]
    Так как
    \[
        \delta_k(x_n) \rightarrow \cos k,
    \]
    то если последовательность сходится, то она сходится к последовательности $x\colon x(k) = \cos k$.

    Таким образом $y_n = x_n - x$ должна слабо сходиться к нулю.
    Пусть $y \in \ell^{\infty}$ определяется как
    \[
        y(k) = \sin k - \cos k = \sqrt{2}\sin\left(k - \frac{\pi}{4}\right).
    \]
    Заметим, что $y_n - y$ -- финитна и, следовательно, при всяком $n \in \N$ лежит в $c_{0}$.
    Рассмотрим $L = \ell^\infty/c_0$ и $\pi\colon \ell^{\infty} \to L$ -- каноническую проекцию.
    Тогда $\pi(y_n) = \pi(y)$.
    С другой стороны, само $y$ не стремится к нулю на бесконечности, а значит $\pi(y) \neq 0$.
    Тогда существует (в силу следствия из теоремы Хана-Банаха) функционал $\Lambda$ в $L^*$ такой, что
    \[
        \Lambda(\pi(y)) = 1.
    \]
    Рассмотрим $\tilde{\Lambda} = \Lambda \circ \pi \in (\ell^\infty)^*$.
    Тогда
    \[
        \tilde{\Lambda}(y_n) = \Lambda(\pi(y_n)) = \Lambda(\pi(y)) = 1 \not\rightarrow 0, n \rightarrow \infty.
    \]
    А значит слабой сходимости последовательности $y_n$ к $0$ нет, и, таким образом, $\{x_n\}$ не является слабо сходящейся.
\end{solution}
\begin{task}[T7]
    Привести пример функционала $f \in (\ell^\infty)^*$, не достигающего своей нормы.
\end{task}
\begin{solution}
    Рассмотрим $c \subset \ell^{\infty}$ и определим на нём функционал $f\colon x \mapsto \lim_{n} x(n)$.
    Он линеен и $|f(x)| = \lim |x(n)| \leq \sup_{n} |x(n)| = \|x\|$ -- значит $\|f\| \leq 1$.
    С другой стороны на векторе $x = (1, \ldots, 1)$ норма достигается, а значит $\|f\| = 1$.

    Продолжим его с сохранением нормы на всё $(\ell^\infty)^*$ и определим $g\colon \ell^\infty \to \R$ как
    \[
        g(x) = \sum_{n = 1}^{\infty} 2^{-n} x(n) - f(x).
    \]
    Очевидно, что $g \in (\ell^{\infty})^*$ и $\|f\| \leq 2$.
    Покажем, что его норма действительно равна $2$.
    Пусть $x^m$ определяется как
    \[
        x^m(n) = \begin{cases}
                     1, & n \leq m, \\
                     -1, & n > m.
        \end{cases}
    \]
    Тогда
    \[
        g(x^m) = \sum_{n = 1}^{m} 2^{-n} - \sum_{n = m + 1}^{\infty} 2^{-n} + 1 = 2 - 2^{1 - m} \rightarrow 2.
    \]
    Пусть норма $g$ достигается на $x\colon \|x\| \leq 1$.
    Так как
    \[
        -1 \leq \sum_{n = 1}^{\infty} 2^{-n}x(n) \leq 1, \quad -1 \leq -f(x) \leq 1,
    \]
    то должно быть одновременно
    \[
        \sum_{n = 1}^{\infty} 2^{-n}x(n) = 1, \quad -f(x) = 1,
    \]
    или то же самое, но с $-1$.
    Но это невозможно: из первого равенства следует, что $x(n) = 1$ при всяком $n$, но тогда $-f(x) = -1$ (если рассматривать равенство $-1$, то $x(n) = -1$, а $-f(x) = 1$).
    Значит $g$ не достигает своей нормы.
\end{solution}
\begin{task}[T8]
    Привести пример несепарабельного банахова пространства $X$ такого, что замкнутый единичный шар в его сопряжённом пространстве $X^*$ не является слабо-* секвенциально компактным.
\end{task}
\begin{solution}
    Рассмотрим $\ell^\infty$.
    Для каждого $n \in \N$ рассмотрим координатные функционалы $\delta_n\colon \ell^\infty \to \R\colon$
    \[
        \delta_n x = x(n).
    \]
    Очевидно, что $\|\delta_n\| = 1$ и все они лежат на единичной сфере $(\ell^\infty)^*$.

    Пусть $\{\delta_{n_k}\}$ -- *-слабо сходящаяся подпоследовательность к некоторому $f \in (\ell^\infty)^*$
    Тогда для всякого $x \in \ell^{\infty}$
    \[
        \delta_{n_k} x \rightarrow f(x),
    \]
    то есть
    \[
        x_{n_k} \rightarrow f(x).
    \]
    Возьмём $x \in \ell^{\infty}$, определяемый как
    \[
        x(n) = \begin{cases}
        (-1)^{k}, & \exists k\colon n = n_k, \\
                   0, & \not\exists k \colon n = n_k.
        \end{cases}
    \]
    Тогда мы получаем, что
    \[
        (-1)^{k} \rightarrow f(x)
    \]
    Чего быть не может, а значит единичная сфера в сопряжённом к $\ell^\infty$ не является *-слабо секвенциально компактной.
\end{solution}
\begin{task}[T9]
    Доказать, что в рефлексивном банаховом пространстве $X$ замкнутый единичный шар $B$ является слабо секвенциально компактным.
\end{task}
\begin{solution}
    В *-слабой топологии на $X^{**}$ единичный шар *-слабо компактен по теореме Банаха-Алаоглу.
    Заметим, что слабая топология на $X$ и *-слабая топология на $X^{**}$ совпадают, а каноническая изометрия Банаха осуществляет гомеоморфизм между $(X, \tau_w)$ и $(X^{**}, \tau_{*w})$.
    Следовательно, единичный шар слабо-компактен в $X$.
    Но, тогда, он и слабо секвенциально компактен по теореме Эберлейна-Шмульяна, что и требовалось показать.
\end{solution}
\begin{task}[T10]
    Пусть в нормированном пространстве $X$ замкнутый единичный шар $B_1[0]$ является слабо компактным.
    Доказать, что $X$ -- рефлексивно.
\end{task}
\begin{solution}
    Пусть $B_1[0]$ является слабо-компактным.
    Пусть $\Phi\colon X \to X^{**}$ -- каноническая изометрия Банаха.
    Для доказательства утверждения задачи достаточно показать сюръективность $\Phi$.

    Заметим, что $\Phi(B_1[0])$ тогда *-слабо компактен в $X^{**}$.
    Так как компакт в хаусдорфовом пространстве замкнут, то $\Phi(B_1[0])$ -- *-слабо замкнут в $X^{**}$.
    В тоже время, согласно задаче T11 он *-слабо плотен в $B \subset X^{**}$ -- единичном шаре дважды сопряжённого пространства.
    Тогда он совпадает с ним.
    Следовательно, $\Phi(X) = X^{**}$ в силу линейности $\Phi$.
\end{solution}
\begin{task}[T11]
    Пусть $X$ -- нормированное пространство, $F\colon X \to X^{**}$ -- изометрия Банаха.
    Доказать, что $F(X)$ является *-слабо всюду плотным в $X^{**}$.
\end{task}
\begin{solution}
    Пусть $L = [F(X)]^{*w}$ и $L \neq X^{**}$.
    Тогда существует $x^{**} \in X^{**} \setminus L$.
    Существует некоторая *-слабая окрестность $U(x^{**})$ такая, что $U(x) \cap L = \emptyset$.
    В неё вложено некоторое множество из базы *-слабой топологии, которое, без ограничения общности имеет вид
    \[
        V = \bigcap_{n = 1}^{N} V(f_n, x^{**}, \epsilon).
    \]
    Рассмотрим функционал $p\colon X^{**} \to \R$, определённый следующим образом:
    \[
        p\colon x^{**} \mapsto \frac{1}{\epsilon} \max_{n} |x^{**}(f_n)|.
    \]
    Так как $V \subset U $, то $V \cap L = \emptyset$, а значит $\forall v \in L$ верно, что $p(v - x^{**}) \geq 1$.

    Теперь пусть $M = L \oplus \Lin\{x^{**}\}$.
    Определим функционал $g\colon M \to \R$ как
    \[
        g\colon v + tx^{**} \mapsto t.
    \]
    Он мажорируется $p\colon$
    \[
        p(v + tx^{**}) = |t|p\left(x^{**} + \frac{v}{t}\right) \geq |t| = |g(tx^{**})| = |g(v + tx^{**})|.
    \]
    Продолжим $g$ на всё $X^{**}$ c сохранением мажоранты по теореме Хана-Банаха.
    Пусть $G\colon X^{**} \to \R$ -- искомое продолжение.
    По определению $G|_{L} = 0$ и $G(x^{**}) = 1$.
    Заметим, что $G$ является *-слабо непрерывным функционалом.
    Тогда существует некоторый $f \in X^{*}$ такой, что $G(x^{**}) = x^{**}(f)$.
    Но $G$ равен $0$ на $L$, а значит и на $F(X)$ и тогда $f = 0$.
    Противоречие, $G(x^{**}) = 1$.

\end{solution}
\begin{task}[T12] % Не нужно применять голдстайна?
    Пусть $X$ -- банахово пространство.
    Доказать, что $X$ -- рефлексивно тогда и только тогда, когда рефлексивно $X^*$.
\end{task}
\begin{solution}
    Пусть $X$ -- рефлексивно.
    Тогда каноническая изометрия Банаха $J_X\colon x \mapsto \Phi_x$, где
    \[
        \Phi_x f = f(x),
    \]
    сюрьективна.
    Пусть $F \in X^{***}$.
    Определим функционал $g \in X^*$ как
    \[
        g(x) = F\left(J_X x\right).
    \]
    Покажем, что если $J_{X^*}$ -- вложение Банаха между $X^*$ и $X^{***}$, то
    \[
        J_{X^*}g = F.
    \]
    Пусть $x^{**} \in X^{**}$.
    Тогда существует $x \in X$ такой, что $J_X x = x^{**}$.
    А значит
    \[
        J_{X^*}g(x^{**}) = x^{**}(g) = J_X x (g) = g(x) = F(J_X x) = F(x^{**}).
    \]
    Таким образом $J_{X^*}$ -- сюръекция, что и показывает рефлексивность $X^*$.

    Пусть теперь $X^*$ -- рефлексивно.
    Рассмотрим оператор $R\colon X^{***} \to X^*$, который действует как
    \[
        \Phi \mapsto \Phi \circ J_X.
    \]
    Пусть $f \in X^*$.
    Положим $\Phi_f = J_{X^*}(f)$.
    Для $x \in X\colon$
    \[
        R(\Phi_f)(x) = \Phi_f(J_X x) = J_{X^*}(f)(J_X x) = J_X x(f) = f(x).
    \]
    А значит $R$ -- сюръективен.
    Пусть $R(\Phi) = 0$.
    Тогда $\Phi = 0$ на $J_X(X)$.
    С другой стороны, $J_X(X)$ является *-слабо плотным в $X^{**}$, а поскольку $X^{**}$ -- рефлексивно, то тогда $\Phi = 0$.
    Таким образом $R$ -- обратный оператор к $J_{X^*}$.

    Заметим, что $R$ -- сопряжён с $J_X$.
    И действительно: если $F \in X^{**}$, то
    \[
        J_X^* F(x) = F(J_X x) = RF(x).
    \]
    Тогда, по теореме Фредгольма
    \[
        \ker R = \left(\im J_X\right)^{\perp}.
    \]
    Но $\ker R = 0$, а значит
    \[
        \left(\im J_X\right)^{\perp} = 0.
    \]
    Навешивая левый аннулятор мы получаем, что
    \[
        [\im J_X] = X^{**}.
    \]
    Но образ $J_X$ и так замкнут, а значит
    \[
        \im J_X = X^{**}.
    \]
    Что и требовалось показать.
\end{solution}
\begin{task}[T13]
    Пусть $X$ -- рефлексивное банахово пространство, $M \subset X$ -- его замкнутое выпуклое подмножество.
    Доказать, что для любого $x \in X$ существует $y \in M$ такой, что $\|x - y \| = \rho(x, M)$.
\end{task}
\begin{solution}
    Пусть $\{y_m\} \subset M$ такова, что
    \[
        \lim_{m \rightarrow \infty} \|x - y_m\| = \rho(x, M).
    \]
    Заметим, что последовательность $z_m = x - y_m$ -- ограничена.
    По теореме Банаха-Тихонова мы можем выделить из неё слабо сходящуюся подпоследовательность $z_{m_k}$ к некоторому $z^*$.
    Заметим, что тогда в силу замкнутости $M$ (согласно теореме Мазура оно и слабо-замкнуто) точка $y^* = x - z^*$ лежит в $M$.
    В силу слабой полунепрерывности нормы снизу и определения инфинума в ней действительно достигается искомое расстояние.
\end{solution}
\begin{task}[T14]
    Пусть $X$ -- рефлексивное банахово пространство, $Y$ -- нормированное и $A \in \LL(X, Y)$.
    Доказать, что $AB_1[0]$ сильно замкнут в $Y$.
\end{task}
\begin{solution}
    Заметим, что в силу выпуклости $B_1[0]$ и линейности $A$ образ шара также является выпуклым.
    Следовательно, по теореме Мазура, для того, чтобы показать сильную замкнутость $AB_1[0]$ достаточно показать его слабую замкнутость.

    По теореме Банаха-Алаоглу (и в силу рефлексивности $X$) $B_1[0]$ является слабо-компактным в $X$.
    Покажем, что $A$ непрерывен как отображение между $X$ и $Y$ со слабыми топологиями.
    Пусть $f \in Y^*$.
    Очевидно, что $f \circ A \in X^*$.
    Заметим, что в силу явного вида базы (то есть пересечение прообразов элементов базы под действием некоторых непрерывных функционалов), всякий элемент базы слабой топологии на $Y$ прообразом переходит в элемент базы слабой топологии на $X$.
    Следовательно, поскольку отображение непрерывно на базе, то оно непрерывно и в целом.
    Тогда $AB_1[0]$ слабо-компактен в $Y$ и, в силу хаусдорфовости слабой топологии, является и слабозамкнутым -- что и требовалось показать.
\end{solution}
\section{Сопряжённые линейные ограниченные операторы.}
\begin{task}[11.1]
    Найти сопряжённый к оператору $A\colon L^2([0, 1]) \to L^2([0, 1])$, если
    \begin{enumerate}
        \item $(Ax)(t) = \int_0^1 t x(s)ds$.
        \item $(Ax)(t) = \int_0^t s x(s)ds$.
    \end{enumerate}
\end{task}
\begin{solution}
    По определению, гильбертов сопряжённый действует как
    \[
        \langle A^* x, y \rangle = \langle x, Ay \rangle, \quad \forall x, y \in L^2([0, 1]).
    \]
    То есть
    \[
        \langle A^{*} x, y \rangle = \int_{0}^{1} x(s) \overline{Ay(s)}ds = \int_{0}^{1}\int_{0}^{1} x(s) s \overline{y(t)}dtds = \int_0^1 \overline{y(t)}dt \cdot \int_0^1 s x(s)ds.
    \]
    Положим
    \[
        A^* x(t) = \int_0^1 s x(s)ds.
    \]
    Тогда
    \[
        \langle A^* x, y \rangle = \int_0^1 A^* x(t) \overline{y(t)}dt = \int_0^1 \int_0^1 s x(s)ds \cdot \overline{y(t)}dt = \int_0^1 s x(s)ds \cdot \int_0^1 \overline{y(t)}dt,
    \]
    что и требовалось.

    Пусть теперь
    \[
        Ax(t) = \int_0^t s x(s) ds.
    \]
    Тогда
    \[
        \langle A^* x, y \rangle = \langle x, Ay \rangle = \int_0^1 x(t) \overline{Ay(t)}dt = \int_0^1 \int_0^t x(t) s \overline{y(s)}dsdt.
    \]
    Используем теорему Фубини:
    \[
        \langle A^* x, y \rangle = \int_0^1 \int_s^{1} x(t)dt \cdot s\overline{y(s)}ds = \int_0^1 s\int_s^1 x(t)dt \cdot \overline{y(s)}ds.
    \]
    То есть
    \[
        A^{*}x(s) = s \int_s^1 x(t)dt.
    \]
\end{solution}
\begin{task}[11.3]
    В пространстве $\ell^2$ для $x \in \ell^2$ положим
    \[
        A_n x = T^n x,
    \]
    где $T$ -- оператор левого сдвига.
    Найти $A_n^*$ и выяснить, являются ли последовательности $A_n$ и $A_n^*$ сходящимися поточечно.
\end{task}
\begin{solution}
    Оператор $T$ действует как
    \[
        Tx(n) = x(n + 1).
    \]
    Гильбертов сопряжённый определяется как
    \[
        \langle T^* x, y \rangle = \langle x, Ty \rangle.
    \]
    Распишем подробнее скалярное произведение:
    \[
        \langle x, Ty \rangle = \sum_{n = 1}^{\infty} x(n) \overline{Ty(n)} = \sum_{n = 1}^{\infty} x(n) \overline{y(n + 1)}.
    \]
    Положим
    \[
        T^* x(n) = \begin{cases}
                       0, & n = 1, \\
                       x(n - 1), & n > 1.
        \end{cases}
    \]
    Тогда
    \[
        \langle T^* x, y \rangle = \sum_{n = 1}^{\infty} T^*x(n) \overline{y}(n) = \sum_{n = 2}^{\infty} x(n - 1)\overline{y(n)} = \sum_{n = 1}^{\infty} x(n)\overline{y(n + 1)}.
    \]
    А значит сопряжённым к оператору левого сдвига является оператор правого сдвига.

    Исследуем поточечную сходимость $T^n$.
    Заметим, что для всякого $x \in \ell^2$
    \[
        \|T^n x\|^2 = \sum_{k = n + 1}^{\infty} |x(k)|^2 \rightarrow 0, n \rightarrow \infty.
    \]
    А значит $\{T^n\}$ поточечно сходится к $0$.

    В тоже время, для $(T^*)^n$ это неверно: $\|(T^*)^n x\| = \|x\|$, но при этом есть покоординатная сходимость к $0$ (а значит, если бы была поточечная сходимость, то сходимость была к 0).
\end{solution}
\begin{task}[11.10]
    Пусть $E$ -- банахово пространство и $A \in \LL(L^2([0, 1]), E)$.
    Пусть $\im A^* \supset C[0, 1]$.
    Найти $\ker A$.
\end{task}
\begin{solution}
    По следствию из теоремы Фредгольма
    \[
        (\ker A)^{\perp} = [\im A^*]^{*-w}.
    \]
    Так как *-слабое замыкание содержит замыкание по норме, а образ сопряжённого оператора содержит $C[0, 1]$, то *-слабое замыкание образа совпадает со всем $L^2([0, 1])$.
    Поскольку в гильбертовых пространствах левый и правый аннуляторы совпадают с ортогональным дополнением, то ортогональное дополнение ядра $A$ совпадает со всем пространством.
    А значит $\ker A = 0$.
\end{solution}
\begin{task}[11.12]
    Пусть $E_1$ и $E_2$ -- нормированные пространства и $A \in \LL(E_1, E_2)$, причём существует $A^{-1} \in \LL(E_2, E_1)$.
    Доказать, что существует $(A^*)^{-1} \in \LL(E_1^*, E_2^*)$, причём $(A^*)^{-1} = (A^{-1})^*$.
\end{task}
\begin{solution}
    Покажем, что для операторов $S\colon E_1 \to E_2$ и $T\colon E_2 \to E_3$ верно следующее:
    \[
        (TS)^* = S^* T^*.
    \]
    Пусть $f \in E_3^*$.
    Тогда $T^* f \in E_2^*$ и
    \[
        (T^* f)(e_2) = f(Te_2).
    \]
    А действие $S^*$ на $T^* f$ имеет вид:
    \[
        S^*(T^* f)(e_1) = (T^* f)(S e_1) = f(T(S e_1)) = f((TS) e_1) = (TS)^{*}f(e_1)
    \]
    по определению.
    Следовательно, оператора $(TS)^*$ и $S^* T^*$ совпадают на всяком векторе, а значит совпадают и в принципе.

    Очевидно, что сопряжённым оператором к $\id_{E_1}$ является $\id_{E_1^*}$ (аналогично и с $E_2$).
    Тогда
    \[
        (A^{-1})^* A^* = (AA^{-1})^* = \id_{E_2}^* = \id_{E_2^*}, \quad A^* (A^{-1})^* = (A^{-1}A)^* = \id_{E_1}^* = \id_{E_1^*}.
    \]
    Следовательно, $(A^{-1})^*$ является обратным к $A^*$.
    В то же время, он ограничен, как сопряжённый к ограниченному, и, таким образом, утверждение задачи показано.
\end{solution}
\begin{task}[T15]
    Доказать, что оператор $A\colon L^2(\R) \to c$ вида
    \[
        (Ax)(k) = \int_{-k}^{k} \frac{x(t)}{1 + t^2}dt, \ \forall k \in \N, \ \forall x \in L^2(\R).
    \]
    является ограниченным и найти его сопряжённый.
    Исследовать множество $AB_1[0]$ на вполне ограниченность и замкнутость в пространстве $c$.
\end{task}
\begin{solution}
    Пусть $g(t) = \frac{1}{1 + t^2}$ и $g_k(t) = \Ind_{[-k, k]}(t)g(t)$.
    Заметим, что $g \in L^2(\R)$.
    Тогда
    \[
        (Ax)(k) = \langle x, g_k \rangle.
    \]
    А значит
    \[
        |(Ax)(k)| \leq \|x\|_2\|g_k\|_2 \leq \|x\|\|g\|.
    \]
    Переходя к супремуму по $k$ получим, что
    \[
        \|Ax\| \leq \|x\|\|g\|.
    \]
    А значит $A$ -- непрерывен и его норма не более чем норма $g$.

    Найдём банахов сопряжённый к $A$.
    Так как $c^* \cong \ell^1$ и $(L^2(\R))^* \cong L^2(\R)$, то банахов сопряжённый действует из $\ell^1$ в $L^2(\R)$.
    По определению, для $f \in \ell^1\colon$
    \[
        A^* f(x) = f(Ax).
    \]
    Функционал действует как
    \begin{multline*}
        f(0)\lim_{k}Ax(k) + \sum_{k = 1}^{\infty} f(k)Ax(k) = f(0) \int_{\R} x(t)g(t)dt + \sum_{k = 1}^{\infty} f(k) \int_{\R}x(t)g_k(t)dt = \\
        = \int_{\R} \frac{1}{1 + t^2}\left(f(0) + \sum_{k = 1}^{\infty}f(k)\Ind_{[-k, k]}\right)dt,
    \end{multline*}
    где перемена суммирования и интегрирования допустима в силу теоремы Фубини.
    Тогда
    \[
        (A^* f)(t) = \frac{1}{1 + t^2} \left(f(0) + \sum_{k = 1}^{\infty}(k)\Ind_{[-k, k]}(t)\right).
    \]

    Покажем, что оператор $A$ -- компактен и является пределом по норме операторов вида
    \[
        A_{N}x = \left(\langle x, g_1 \rangle, \ldots, \langle x, g_N \rangle, \langle x, g \rangle, \ldots   \right).
    \]
    Очевидно, что $\im A_{N} \subset \Lin\{e_0, e_1, e_2, \ldots, e_N\}$ и образ $A_N$ -- конечномерен.
    В тоже время
    \[
        \|(A - A_N)x\| = \sup_{k > N} |\langle x, g_k - g \rangle| \leq \|x\|\sup_{k > N} \|g_k - g\| \leq \|x\|\|g_N - g\|.
    \]
    Из чего и следует сходимость по норме.
    Тогда $A$ компактен как равномерный предел конечномерных операторов, а значит он переводит шар $B_1[0]$ во вполне ограниченное множество.

    Пусть $y^{(n)} = Ax_n \in AB_1[0]$ и $y^{(n)} \rightarrow y \in c$.
    Так как $\{x_n\}$ -- ограничена, то из неё можно выделить слабо сходящуюся к некоторому $x$ подпоследовательность $\{x_{n_k}\}$.
    В силу слабой непрерывности $x \mapsto (Ax)(k)$ верно, что $(Ax_{n_j})(k) \rightarrow Ax(k)$.
    С другой стороны, так как из сходимости по $c$-норме следует покоординатная, то $(Ax_{n_j})(k) \rightarrow y(k)$, и $y = Ax \in AB_1[0]$ -- а значит $AB_1[0]$ замкнуто.
\end{solution}
\begin{task}[T16]
    Доказать, что оператор $A\colon \ell^1 \to L^1((0, +\infty))$ вида
    \[
        (Ax)(t) = \sum_{k = 1}^{\infty} x(k)\exp(-kt), \ \text{a.s. } t > 0, \ \forall x \in \ell^1,
    \]
    является ограниченным, и найти его сопряжённый.
    Исследовать множество $AB_1[0]$ на вполне ограниченность и замкнутость в пространстве $L^1((0, +\infty))$.
\end{task}
\begin{solution}
    Исследуем оператор на ограниченность.
    Пусть $x \in \ell^1$.
    Тогда
    \begin{multline*}
        \|Ax\|_1 = \int_{\R_{++}} |Ax(t)|dt = \int_{\R_++} \left|\sum_{k = 1}^{\infty}x(k)\exp(-kt)\right| \leq \int_{\R_{++}}\sum_{k = 1}^{\infty} |x(k)|\exp(-kt)dt = \\ = \sum_{k = 1}^{\infty}|x(k)| \int_{\R_{++}} \exp(-kt)dt = \sum_{k = 1}^{\infty} \frac{|x(k)|}{k} \leq \sum_{k = 1}^{\infty} |x(k)| = \|x\|_1.
    \end{multline*}
    Что и показывает ограниченность оператора.

    Найдём сопряжённый $A^*\colon L^\infty((0, +\infty)) \to \ell^{\infty}$.
    Пусть $f \in L^{\infty}((0, +\infty))$ и $x \in \ell^1$.
    По определению банахова сопряжённого:
    \[
        (A^{*}f)(x) = f(Ax).
    \]
    То есть
    \begin{multline*}
    (A^{*}f)(x) = \int_{\R_{++}} f(t)Ax(t)dt = \int_{\R_{++}} f(t) \sum_{k = 1}^{\infty} x(k)\exp(-kt)dt = \\
    = \sum_{k = 1}^{\infty} x(k) \int_{\R_{++}} f(t) \exp(-kt) dt.
    \end{multline*}
    И тогда
    \[
        (A^*f)(k) = \int_{\R_{++}} f(t)\exp(-kt)dt.
    \]

    Рассмотрим срезки оператора $A\colon$
    \[
        (A_N x)(t) = \sum_{k = 1}^{N} x(k)\exp(-kt).
    \]
    Образ $A_N$ -- конечномерен, а
    \[
        \|(A - A_N)x\| = \int_0^\infty \left|\sum_{k = N + 1}^{\infty} x(k)\exp(-kt)\right| \leq \sum_{k = N + 1}^{\infty} \frac{|x(k)|}{k} \leq \frac{\|x\|}{N + 1}.
    \]
    А значит у нас оборудована сходимость $A_N$ по норме к $A$.
    Следовательно, $A$ -- компактен и переводит шар во вполне ограниченное множество.

    Покажем замкнутость образа $B_1[0]$.
    Пусть $y_n = Ax_n \subset AB_1[0]$ и $y_n \rightarrow y \in L^1((0, +\infty))$.
    Так как $\forall k \in \N$ $|x_n(k)| \leq \|x_n\| \leq 1$, то все координатные последовательности ограничены и допускают выделение фундаментальной подпоследовательности.
    Выберем подпоследовательность $\{x_{n_j}\}$ такую, что
    \[
        \forall k \in \N \hookrightarrow x_{n_j}(k) \rightarrow x(k),
    \]
    для некоторого $x \in \ell^1$.
    В силу леммы Фату $x \in B_1[0]$.
    Покажем сходимость $Ax_{n_j}$ к $Ax$.
    Для всякого $N \in \N\colon$
    \[
        \|Ax_{n_j} - Ax\| \leq \sum_{k = 1}^{N} \frac{|x_{n_j}(k) - x(k)|}{k} + \sum_{k = N + 1} \frac{|x_{n_j}(k)| + |x(k)|}{k} \leq \sum_{k = 1}^{N} \frac{|x_{n_j}(k) - x(k)|}{k} + \frac{2}{N + 1}.
    \]
    Первая сумма стремится к нулю при фиксированном $N \in \N$, второе слагаемое можно бесконечно малым в силу произвольности выбора $N$.
    Тогда у нас есть искомая сходимость и значит $y = Ax$, то есть $AB_1[0]$ -- замкнуто.
\end{solution}
